% Generated draft: requires mathematical and source review.
\documentclass[11pt]{article}
\usepackage[margin=25mm]{geometry}
\usepackage{fontspec}
\setmainfont{Latin Modern Roman}
\usepackage{amsmath,amssymb,mathtools,unicode-math}
\setmathfont{Latin Modern Math}
\usepackage{enumitem,needspace,xurl,hyperref}
\setlength{\parindent}{0pt}
\setlength{\parskip}{5pt}
\setlength{\emergencystretch}{3em}
\newcommand{\N}{\mathbb{N}}
\newcommand{\Z}{\mathbb{Z}}
\newcommand{\Q}{\mathbb{Q}}
\newcommand{\R}{\mathbb{R}}
\newcommand{\C}{\mathbb{C}}
\newcommand{\F}{\mathbb{F}}
\newcommand{\A}{\mathbb{A}}
\newcommand{\PP}{\mathbb{P}}
\newcommand{\E}{\mathbb{E}}
\newcommand{\eps}{\varepsilon}
\newcommand{\dd}{\,\mathrm{d}}
\begin{document}
\section*{UnsolvedMath ID: 801149}
\textbf{Source problem code:} catalogue entry
\subsection{Definitions and mathematical statement}
For a tangent-bundle endomorphism \(L:TM\to TM\), its Nijenhuis torsion is the vector-valued two-form defined, for vector fields \(X,Y\), by
\[
\mathcal N_L(X,Y)=[LX,LY]-L[LX,Y]-L[X,LY]+L^2[X,Y].
\]
The operator \(L\) is called Nijenhuis when \(\mathcal N_L=0\). The conventional pointwise meaning of \(\mathrm{gl}\)-regularity on an \(n\)-dimensional manifold is that, at each point \(p\), the minimal polynomial of \(L_p\) has degree \(n\); equivalently, the centralizer of \(L_p\) in \(\operatorname{End}(T_pM)\) has dimension \(n\). These definitions explain the terminology but do not supply the missing geometric or regularity hypotheses of the source excerpt.

\textbf{Problem (supplied formulation, incomplete context).} The source refers to a system of equations involving \(L\) and auxiliary data denoted by \(e\) and \(\alpha\), but does not provide that system in the supplied record. Can this system be written so that its coefficients depend only on \(L\), at least when \(L\) is \(\mathrm{gl}\)-regular? Are its compatibility conditions precisely
\[
\mathcal N_L=0?
\]
The source further requires these compatibility conditions to be geometric conditions on \(L\) that vanish whenever \(L\) is a Nijenhuis operator.

As motivation, the source states that, in the \(\mathrm{gl}\)-regular case, any two pairs \((L_1,e_1)\) and \((L_2,e_2)\) with identical Jordan blocks of \(L_1\) and \(L_2\) are isomorphic by a coordinate transformation, and that \(\alpha\) is the dual form for \(e\). These are retained as source assertions whose hypotheses and meanings require the omitted context.

\textbf{Unresolved qualifications.} The equations, their unknowns, the definitions and constraints on \(e\) and \(\alpha\), and the intended meaning of compatibility are absent. The manifold's dimension, regularity category, local or global scope, and any boundary or initial conditions are also unspecified. Consequently, the excerpt determines two questions about an unidentified system, but does not determine a fully quantified mathematical statement or a precise equivalence between compatibility and vanishing Nijenhuis torsion.
\subsection{Short English statement}
The problem asks whether an unspecified system can be expressed using coefficients determined solely by \(L\), at least in the \(\mathrm{gl}\)-regular case, and whether its compatibility conditions are exactly the vanishing of the Nijenhuis torsion. The supplied excerpt omits the system and its auxiliary definitions, so a precise formulation requires additional source context. The registry labels the problem open; that label has not been independently verified.
\subsection{Sources}
\begin{itemize}
\item Local UnsolvedMath registry, permanent ID \(801149\), exact source problem code \texttt{catalogue entry}, title \emph{Nijenhuis Geometry and Integrable Systems --- Problem 2}. The supplied record attributes the problem to Bolsinov, Konyaev, and Matveev, with proposed year 2024, and labels it open.
\item Bolsinov, Konyaev, Matveev, \emph{Research Problems on Nijenhuis Geometry and Integrable Systems} (2024), Problem 2. Supplied source URL: \url{https://arxiv.org/abs/2410.04276}. The provenance records access on 2026-10-16 and extraction by \texttt{arxiv\_latex}; these are supplied provenance details, not a claim of independent source access or verification.
\end{itemize}
\end{document}
